\label{app:robustorder}

Fix a set of systems $S$, a set of items, and a grid $C$ of evaluation cells. In cell $c$ each
system has a mean score, and those means induce a total order $\prec_c$ on $S$. A published table
reports $\prec_{c_0}$ for one cell $c_0$ that its paper does not name. Write
\[
  \prec_{\!*} \;=\; \bigcap_{c\in C} \prec_c ,
\]
and say $a$ \emph{dominates} $b$ when $a \prec_{\!*} b$. Say the pair is \emph{separated in every
cell} when in addition the paired bootstrap interval on the margin excludes zero in each $c$.

\begin{quote}
\textbf{Proposition 1.} $\prec_{\!*}$ is a strict partial order, and it is the finest relation on
$S$ that no cell of $C$ can contradict.

\emph{Proof.} An intersection of transitive, irreflexive relations is transitive and irreflexive.
Any relation not contained in some $\prec_c$ contains a pair that cell orders the other way.
\hfill$\square$
\end{quote}

\begin{quote}
\textbf{Proposition 2.} If $C \subseteq C'$ then $\prec_{\!*}^{C'} \subseteq \prec_{\!*}^{C}$: the
dominating set is monotone non-increasing in the grid, and so is the number of tiers.

\emph{Proof.} Immediate from the definition as an intersection; the tier count is the longest chain
of a smaller relation. \hfill$\square$
\end{quote}

Proposition 2 is the reason a single share is not a property of a leaderboard. It is a property of
the pair (leaderboard, grid), it can only fall as a grid is widened, and any paper reporting one
number could have reported a larger one by declaring fewer choices. The remedy used here is to
report the curve rather than a point: each declared axis alone, and their product. The tables below
give that curve for every leaderboard in the paper, together with the accuracy of every system in
every cell, which is the table all the derived quantities are computed from.

\paragraph{Three verdicts, not two.} A pair that does not dominate has failed in one of two ways
that should not be added together. Some cell may put the two systems the other way round with a
paired interval excluding zero, which is a choice deciding against the published ordering; or no
cell may do so, in which case the published ordering merely fails to hold everywhere on a margin
the benchmark cannot resolve. The first is called \emph{contested} and the second
\emph{unresolved}, and only the count of contested pairs is evidence that a choice reverses
anything. Counting sign changes alone would let noise on a tied pair read as a reversal, which is
the error this instrument exists to avoid.

\paragraph{A dual statistic.} The tier count says what survives; the number of distinct total
orders the same systems take across the grid says how much varies. Neither is recoverable from the
other: a leaderboard can take many orderings and still support several tiers if the systems that
move are adjacent, and it can take few and support one if the pairs that never move are also never
resolved.

\paragraph{The grid is the same shape and not the same cells.} The criterion axis is the four
identity relations in every domain. The budget axis is not transferable: retrosynthesis tables are
read at top-1 through top-10 and metabolite work at fifteen, so each leaderboard is swept over the
budgets its own field reports at. Fixing one set of budgets across domains would have measured
robustness to a budget nobody uses. What this costs is that the three shares are comparable in
construction rather than in cells, which is why they are reported as three numbers and not averaged.
One consequence is worth stating because it bounds how much the grid can find: on the seven-system
board exactly one of the $120$ pairs of cells gives identical hit vectors for every system,
canonical and InChIKey at top-1, so the sixteen cells are sixteen and not fewer.

\paragraph{Which leaderboards the grid applies to, decided before the shares were read.} A grid
cell has to be a choice the evaluation could have been made under. Two conditions decide it, and
both exclude cells rather than admit them. First, a matching criterion has to be an identity
relation: \textsc{tanimoto1} is reported in Table~\ref{tab:matchsens}, because published work
compares Morgan fingerprints, but it is not a grid cell, since at radius $2$ over $2048$ bits it
returns $1$ for decane against undecane and for D- against L-alanine, and a rule that identifies a
homologue with its reference is not a candidate answer to whether two structures are the same
compound. Second, the metric must be defined independently of the criterion. That is what excludes
the molecular-generation set: its headline quantity counts \emph{distinct} canonical strings, so the
criterion is not a parameter of the evaluation but part of the definition of what is being counted,
and re-scoring it under another criterion changes the metric rather than the convention. The
generation results are therefore reported in Appendix~\ref{app:xdomain} as a criterion sensitivity
and not as a row here. Both rules were fixed before the shares below were read, and the second
excludes the one leaderboard whose pairs are furthest apart, which would have raised the shares.

\paragraph{What the share does not measure.} It is not an error rate and no cell is privileged.
The instrument takes the grid as declared and says what part of a published ordering is invariant
to it; whether a given cell is one a field would defend is a question about the field, argued in
the text and not settled by the instrument. Nor is the interval on a share a statement about the
grid: it resamples items with the grid held fixed, so it reports how much of the share is the
sample. The grid's contribution is what Proposition 2 and the curve below report.

% generated by scripts/make_robust_tables.py -- do not edit by hand

\paragraph{The share as a function of the grid.} Domination is an intersection over cells, so adding a cell can only remove pairs from it and the share is a property of the pair (leaderboard, grid). Each declared axis is therefore reported alone and against the product.

\begin{center}\small
\begin{tabular}{llcccc}
\toprule
leaderboard & sub-grid & cells & dominate & tiers & orderings \\
\midrule
retrosynthesis, 7 systems & criteria only, at the published budget & 4 & 16 of 21 & 5 & 2 \\
 & budgets only, at the published criterion & 4 & 12 of 21 & 4 & 3 \\
 & the product & 16 & 8 of 21 & 3 & 4 \\
\midrule
retrosynthesis, 3 systems & criteria only, at the published budget & 4 & 3 of 3 & 3 & 1 \\
 & budgets only, at the published criterion & 4 & 0 of 3 & 1 & 3 \\
 & the product & 16 & 0 of 3 & 1 & 3 \\
\midrule
metabolites, 3 methods & criteria only, at the published budget & 4 & 3 of 3 & 3 & 1 \\
 & budgets only, at the published criterion & 4 & 2 of 3 & 2 & 2 \\
 & the product & 16 & 2 of 3 & 2 & 2 \\
\midrule
translation, WMT24 en-de, 19 systems & criteria only, unfloored & 4 & 141 of 171 & 11 & 4 \\
 & the floor only, at the published criterion & 2 & 159 of 171 & 15 & 2 \\
 & the product & 8 & 135 of 171 & 11 & 6 \\
\midrule
translation, WMT24 ja-zh, 15 systems & criteria only, unfloored & 4 & 97 of 105 & 9 & 3 \\
 & the floor only, at the published criterion & 2 & 102 of 105 & 12 & 2 \\
 & the product & 8 & 97 of 105 & 9 & 5 \\
\midrule
docking, 7 programs & criteria only, at the published post-processing & 4 & 18 of 21 & 5 & 3 \\
 & post-processing only, at the published criterion & 2 & 18 of 21 & 5 & 2 \\
 & the product & 8 & 17 of 21 & 4 & 5 \\
\bottomrule\end{tabular}\end{center}

\paragraph{retrosynthesis, 7 systems.} Published order at ('canonical', 1): graphretro $>$ localretro $>$ retroformer $>$ graph2smiles $>$ gln $>$ g2retro $>$ retroxpert. 4999 items. Accuracy by system and cell, from which every quantity below follows:

\begin{center}\small
\begin{tabular}{lcccccccc}
\toprule
system & canon@1 & canon@3 & canon@5 & canon@10 & no-st@1 & no-st@3 & no-st@5 & no-st@10 \\
\midrule
graphretro & 0.538 & 0.677 & 0.716 & 0.745 & 0.550 & 0.693 & 0.733 & 0.762 \\
localretro & 0.534 & 0.769 & 0.842 & 0.910 & 0.547 & 0.788 & 0.862 & 0.930 \\
retroformer & 0.532 & 0.706 & 0.761 & 0.819 & 0.554 & 0.722 & 0.776 & 0.832 \\
graph2smiles & 0.529 & 0.661 & 0.695 & 0.724 & 0.557 & 0.679 & 0.711 & 0.738 \\
gln & 0.525 & 0.690 & 0.756 & 0.837 & 0.538 & 0.702 & 0.768 & 0.849 \\
g2retro & 0.515 & 0.722 & 0.783 & 0.837 & 0.529 & 0.741 & 0.803 & 0.858 \\
retroxpert & 0.447 & 0.598 & 0.644 & 0.694 & 0.462 & 0.610 & 0.655 & 0.703 \\
\bottomrule\end{tabular}\end{center}
\begin{center}\small
\begin{tabular}{lcccccccc}
\toprule
system & IK@1 & IK@3 & IK@5 & IK@10 & taut@1 & taut@3 & taut@5 & taut@10 \\
\midrule
graphretro & 0.538 & 0.677 & 0.716 & 0.745 & 0.555 & 0.701 & 0.741 & 0.770 \\
localretro & 0.534 & 0.769 & 0.842 & 0.910 & 0.549 & 0.791 & 0.865 & 0.932 \\
retroformer & 0.532 & 0.706 & 0.762 & 0.819 & 0.556 & 0.724 & 0.778 & 0.834 \\
graph2smiles & 0.529 & 0.661 & 0.695 & 0.724 & 0.559 & 0.682 & 0.714 & 0.740 \\
gln & 0.525 & 0.690 & 0.756 & 0.837 & 0.540 & 0.705 & 0.771 & 0.852 \\
g2retro & 0.515 & 0.722 & 0.783 & 0.837 & 0.531 & 0.743 & 0.806 & 0.862 \\
retroxpert & 0.447 & 0.598 & 0.644 & 0.694 & 0.463 & 0.612 & 0.657 & 0.705 \\
\bottomrule\end{tabular}\end{center}

Per-pair classification. A pair is \emph{contested} only when some cell puts it the other way round with an interval excluding zero; a pair that is neither dominating nor contested is one the benchmark does not resolve, which is a different failure and is not counted as a reversal.

\begin{center}\small
\begin{tabular}{llccccl}
\toprule
pair (as the published cell orders it) & verdict & sep.\ every cell & cells reversing & with an interval & own cell resolves \\
\midrule
gln over g2retro & contested & no & 12 & 8 & no \\
graph2smiles over g2retro & contested & no & 12 & 12 & yes \\
graphretro over g2retro & contested & no & 12 & 12 & yes \\
localretro over g2retro & dominates & yes & 0 & 0 & yes \\
retroformer over g2retro & contested & no & 12 & 12 & yes \\
g2retro over retroxpert & dominates & yes & 0 & 0 & yes \\
graph2smiles over gln & contested & no & 12 & 12 & no \\
graphretro over gln & contested & no & 12 & 8 & no \\
localretro over gln & dominates & no & 0 & 0 & no \\
retroformer over gln & contested & no & 4 & 4 & no \\
gln over retroxpert & dominates & yes & 0 & 0 & yes \\
graphretro over graph2smiles & unresolved & no & 2 & 0 & no \\
localretro over graph2smiles & unresolved & no & 2 & 0 & no \\
retroformer over graph2smiles & unresolved & no & 2 & 0 & no \\
graph2smiles over retroxpert & dominates & yes & 0 & 0 & yes \\
graphretro over localretro & contested & no & 12 & 12 & no \\
graphretro over retroformer & contested & no & 14 & 12 & no \\
graphretro over retroxpert & dominates & yes & 0 & 0 & yes \\
localretro over retroformer & unresolved & no & 2 & 0 & no \\
localretro over retroxpert & dominates & yes & 0 & 0 & yes \\
retroformer over retroxpert & dominates & yes & 0 & 0 & yes \\
\bottomrule\end{tabular}\end{center}

\paragraph{retrosynthesis, 3 systems.} Published order at ('canonical', 1): chemformer $>$ gta $>$ tiedtransformer. 4982 items. Accuracy by system and cell, from which every quantity below follows:

\begin{center}\small
\begin{tabular}{lcccccccc}
\toprule
system & canon@1 & canon@3 & canon@5 & canon@10 & no-st@1 & no-st@3 & no-st@5 & no-st@10 \\
\midrule
chemformer & 0.539 & 0.609 & 0.618 & 0.624 & 0.550 & 0.621 & 0.630 & 0.636 \\
gta & 0.518 & 0.678 & 0.748 & 0.818 & 0.540 & 0.697 & 0.767 & 0.832 \\
tiedtransformer & 0.471 & 0.679 & 0.743 & 0.798 & 0.481 & 0.693 & 0.756 & 0.811 \\
\bottomrule\end{tabular}\end{center}
\begin{center}\small
\begin{tabular}{lcccccccc}
\toprule
system & IK@1 & IK@3 & IK@5 & IK@10 & taut@1 & taut@3 & taut@5 & taut@10 \\
\midrule
chemformer & 0.539 & 0.609 & 0.618 & 0.624 & 0.553 & 0.624 & 0.634 & 0.639 \\
gta & 0.518 & 0.678 & 0.749 & 0.818 & 0.544 & 0.702 & 0.771 & 0.835 \\
tiedtransformer & 0.471 & 0.679 & 0.743 & 0.798 & 0.484 & 0.697 & 0.761 & 0.815 \\
\bottomrule\end{tabular}\end{center}

Per-pair classification. A pair is \emph{contested} only when some cell puts it the other way round with an interval excluding zero; a pair that is neither dominating nor contested is one the benchmark does not resolve, which is a different failure and is not counted as a reversal.

\begin{center}\small
\begin{tabular}{llccccl}
\toprule
pair (as the published cell orders it) & verdict & sep.\ every cell & cells reversing & with an interval & own cell resolves \\
\midrule
chemformer over gta & contested & no & 12 & 12 & yes \\
chemformer over tiedtransformer & contested & no & 12 & 12 & yes \\
gta over tiedtransformer & unresolved & no & 2 & 0 & yes \\
\bottomrule\end{tabular}\end{center}

\paragraph{metabolites, 3 methods.} Published order at ('tautomer', 15): MetaPredictor $>$ SyGMa $>$ GRAIL. 1170 items. Accuracy by system and cell, from which every quantity below follows:

\begin{center}\small
\begin{tabular}{lcccccccc}
\toprule
system & canon@1 & canon@5 & canon@10 & canon@15 & IK@1 & IK@5 & IK@10 & IK@15 \\
\midrule
MetaPredictor & 0.332 & 0.562 & 0.649 & 0.672 & 0.342 & 0.585 & 0.673 & 0.697 \\
SyGMa & 0.313 & 0.591 & 0.642 & 0.658 & 0.334 & 0.615 & 0.668 & 0.685 \\
GRAIL & 0.248 & 0.403 & 0.431 & 0.439 & 0.253 & 0.411 & 0.440 & 0.449 \\
\bottomrule\end{tabular}\end{center}
\begin{center}\small
\begin{tabular}{lcccccccc}
\toprule
system & no-st@1 & no-st@5 & no-st@10 & no-st@15 & taut@1 & taut@5 & taut@10 & taut@15 \\
\midrule
MetaPredictor & 0.353 & 0.591 & 0.677 & 0.702 & 0.357 & 0.602 & 0.688 & 0.711 \\
SyGMa & 0.334 & 0.615 & 0.669 & 0.688 & 0.340 & 0.626 & 0.681 & 0.698 \\
GRAIL & 0.255 & 0.413 & 0.442 & 0.450 & 0.259 & 0.420 & 0.450 & 0.459 \\
\bottomrule\end{tabular}\end{center}

Per-pair classification. A pair is \emph{contested} only when some cell puts it the other way round with an interval excluding zero; a pair that is neither dominating nor contested is one the benchmark does not resolve, which is a different failure and is not counted as a reversal.

\begin{center}\small
\begin{tabular}{llccccl}
\toprule
pair (as the published cell orders it) & verdict & sep.\ every cell & cells reversing & with an interval & own cell resolves \\
\midrule
MetaPredictor over GRAIL & dominates & yes & 0 & 0 & yes \\
SyGMa over GRAIL & dominates & yes & 0 & 0 & yes \\
MetaPredictor over SyGMa & contested & no & 4 & 2 & no \\
\bottomrule\end{tabular}\end{center}

\paragraph{translation, WMT24 en-de, 19 systems.} Published order at ('wmt', 'none'): Dubformer $>$ GPT-4 $>$ Unbabel-Tower70B $>$ TranssionMT $>$ ONLINE-B $>$ refB $>$ Mistral-Large $>$ CommandR-plus $>$ refA $>$ Gemini-1.5-Pro $>$ ONLINE-W $>$ Claude-3.5 $>$ IOL-Research $>$ Aya23 $>$ ONLINE-A $>$ Llama3-70B $>$ IKUN $>$ IKUN-C $>$ MSLC. 486 items. Accuracy by system and cell, from which every quantity below follows:

\begin{center}\small
\begin{tabular}{lcccccccc}
\toprule
system & wmt / none & wmt / clipped & plain / none & plain / clipped & major-only / none & major-only / clipped & count / none & count / clipped \\
\midrule
Dubformer & -1.584 & -1.583 & -2.432 & -2.346 & -0.292 & -0.292 & -1.263 & -1.261 \\
GPT-4 & -1.587 & -1.587 & -2.368 & -2.294 & -0.261 & -0.261 & -1.323 & -1.323 \\
Unbabel-Tower70B & -1.649 & -1.625 & -2.249 & -2.066 & -0.292 & -0.292 & -1.080 & -1.080 \\
TranssionMT & -1.813 & -1.813 & -2.397 & -2.286 & -0.327 & -0.327 & -1.089 & -1.089 \\
ONLINE-B & -1.817 & -1.807 & -2.360 & -2.243 & -0.315 & -0.315 & -1.101 & -1.101 \\
refB & -1.847 & -1.749 & -2.510 & -2.243 & -0.333 & -0.333 & -1.177 & -1.163 \\
Mistral-Large & -1.930 & -1.930 & -2.584 & -2.523 & -0.333 & -0.333 & -1.251 & -1.251 \\
CommandR-plus & -2.010 & -1.990 & -2.716 & -2.578 & -0.346 & -0.346 & -1.333 & -1.333 \\
refA & -2.121 & -1.997 & -2.753 & -2.545 & -0.377 & -0.377 & -1.247 & -1.232 \\
Gemini-1.5-Pro & -2.204 & -1.834 & -2.296 & -2.216 & -0.307 & -0.307 & -1.070 & -1.070 \\
ONLINE-W & -2.224 & -2.193 & -2.934 & -2.796 & -0.372 & -0.372 & -1.444 & -1.444 \\
Claude-3.5 & -2.281 & -1.812 & -2.763 & -2.401 & -0.354 & -0.354 & -1.348 & -1.348 \\
IOL-Research & -2.392 & -2.351 & -3.136 & -2.965 & -0.401 & -0.401 & -1.531 & -1.531 \\
Aya23 & -3.095 & -3.080 & -3.753 & -3.646 & -0.533 & -0.533 & -1.621 & -1.621 \\
ONLINE-A & -3.304 & -3.273 & -4.078 & -3.922 & -0.568 & -0.568 & -1.807 & -1.807 \\
Llama3-70B & -3.630 & -3.501 & -4.496 & -4.189 & -0.636 & -0.636 & -1.953 & -1.940 \\
IKUN & -3.869 & -3.706 & -4.576 & -4.232 & -0.677 & -0.677 & -1.868 & -1.868 \\
IKUN-C & -5.082 & -4.897 & -5.590 & -5.288 & -0.872 & -0.872 & -2.101 & -2.101 \\
MSLC & -13.483 & -10.940 & -13.412 & -10.726 & -2.442 & -2.442 & -3.642 & -3.634 \\
\bottomrule\end{tabular}\end{center}

Per-pair classification. A pair is \emph{contested} only when some cell puts it the other way round with an interval excluding zero; a pair that is neither dominating nor contested is one the benchmark does not resolve, which is a different failure and is not counted as a reversal.

\begin{center}\small
\begin{tabular}{llccccl}
\toprule
pair (as the published cell orders it) & verdict & sep.\ every cell & cells reversing & with an interval & own cell resolves \\
\midrule
Claude-3.5 over Aya23 & dominates & yes & 0 & 0 & yes \\
CommandR-plus over Aya23 & dominates & yes & 0 & 0 & yes \\
Dubformer over Aya23 & dominates & yes & 0 & 0 & yes \\
GPT-4 over Aya23 & dominates & yes & 0 & 0 & yes \\
Gemini-1.5-Pro over Aya23 & dominates & yes & 0 & 0 & yes \\
Aya23 over IKUN & dominates & yes & 0 & 0 & yes \\
Aya23 over IKUN-C & dominates & yes & 0 & 0 & yes \\
IOL-Research over Aya23 & dominates & no & 0 & 0 & yes \\
Aya23 over Llama3-70B & dominates & no & 0 & 0 & yes \\
Aya23 over MSLC & dominates & yes & 0 & 0 & yes \\
Mistral-Large over Aya23 & dominates & yes & 0 & 0 & yes \\
Aya23 over ONLINE-A & dominates & no & 0 & 0 & no \\
ONLINE-B over Aya23 & dominates & yes & 0 & 0 & yes \\
ONLINE-W over Aya23 & dominates & yes & 0 & 0 & yes \\
TranssionMT over Aya23 & dominates & yes & 0 & 0 & yes \\
Unbabel-Tower70B over Aya23 & dominates & yes & 0 & 0 & yes \\
refA over Aya23 & dominates & yes & 0 & 0 & yes \\
refB over Aya23 & dominates & yes & 0 & 0 & yes \\
CommandR-plus over Claude-3.5 & unresolved & no & 2 & 0 & no \\
Dubformer over Claude-3.5 & dominates & no & 0 & 0 & yes \\
GPT-4 over Claude-3.5 & dominates & no & 0 & 0 & yes \\
Gemini-1.5-Pro over Claude-3.5 & unresolved & no & 1 & 0 & no \\
Claude-3.5 over IKUN & dominates & yes & 0 & 0 & yes \\
Claude-3.5 over IKUN-C & dominates & yes & 0 & 0 & yes \\
Claude-3.5 over IOL-Research & dominates & no & 0 & 0 & no \\
Claude-3.5 over Llama3-70B & dominates & yes & 0 & 0 & yes \\
Claude-3.5 over MSLC & dominates & yes & 0 & 0 & yes \\
Mistral-Large over Claude-3.5 & unresolved & no & 2 & 0 & no \\
Claude-3.5 over ONLINE-A & dominates & yes & 0 & 0 & yes \\
ONLINE-B over Claude-3.5 & dominates & no & 0 & 0 & no \\
ONLINE-W over Claude-3.5 & contested & no & 7 & 2 & no \\
TranssionMT over Claude-3.5 & unresolved & no & 1 & 0 & no \\
Unbabel-Tower70B over Claude-3.5 & dominates & no & 0 & 0 & yes \\
refA over Claude-3.5 & unresolved & no & 4 & 0 & no \\
refB over Claude-3.5 & dominates & no & 0 & 0 & no \\
Dubformer over CommandR-plus & dominates & no & 0 & 0 & yes \\
GPT-4 over CommandR-plus & dominates & no & 0 & 0 & yes \\
CommandR-plus over Gemini-1.5-Pro & contested & no & 7 & 2 & no \\
CommandR-plus over IKUN & dominates & yes & 0 & 0 & yes \\
CommandR-plus over IKUN-C & dominates & yes & 0 & 0 & yes \\
CommandR-plus over IOL-Research & dominates & no & 0 & 0 & yes \\
CommandR-plus over Llama3-70B & dominates & yes & 0 & 0 & yes \\
CommandR-plus over MSLC & dominates & yes & 0 & 0 & yes \\
Mistral-Large over CommandR-plus & dominates & no & 0 & 0 & no \\
CommandR-plus over ONLINE-A & dominates & yes & 0 & 0 & yes \\
ONLINE-B over CommandR-plus & dominates & no & 0 & 0 & no \\
CommandR-plus over ONLINE-W & dominates & no & 0 & 0 & no \\
TranssionMT over CommandR-plus & dominates & no & 0 & 0 & no \\
Unbabel-Tower70B over CommandR-plus & dominates & no & 0 & 0 & yes \\
CommandR-plus over refA & unresolved & no & 3 & 0 & no \\
refB over CommandR-plus & dominates & no & 0 & 0 & no \\
Dubformer over GPT-4 & unresolved & no & 4 & 0 & no \\
Dubformer over Gemini-1.5-Pro & contested & no & 4 & 2 & no \\
Dubformer over IKUN & dominates & yes & 0 & 0 & yes \\
Dubformer over IKUN-C & dominates & yes & 0 & 0 & yes \\
Dubformer over IOL-Research & dominates & yes & 0 & 0 & yes \\
Dubformer over Llama3-70B & dominates & yes & 0 & 0 & yes \\
Dubformer over MSLC & dominates & yes & 0 & 0 & yes \\
Dubformer over Mistral-Large & unresolved & no & 2 & 0 & no \\
Dubformer over ONLINE-A & dominates & yes & 0 & 0 & yes \\
Dubformer over ONLINE-B & contested & no & 4 & 2 & no \\
Dubformer over ONLINE-W & dominates & yes & 0 & 0 & yes \\
Dubformer over TranssionMT & contested & no & 4 & 2 & no \\
Dubformer over Unbabel-Tower70B & contested & no & 6 & 2 & no \\
Dubformer over refA & unresolved & no & 2 & 0 & yes \\
Dubformer over refB & unresolved & no & 3 & 0 & no \\
GPT-4 over Gemini-1.5-Pro & contested & no & 4 & 2 & yes \\
GPT-4 over IKUN & dominates & yes & 0 & 0 & yes \\
GPT-4 over IKUN-C & dominates & yes & 0 & 0 & yes \\
GPT-4 over IOL-Research & dominates & yes & 0 & 0 & yes \\
GPT-4 over Llama3-70B & dominates & yes & 0 & 0 & yes \\
GPT-4 over MSLC & dominates & yes & 0 & 0 & yes \\
GPT-4 over Mistral-Large & unresolved & no & 2 & 0 & yes \\
GPT-4 over ONLINE-A & dominates & yes & 0 & 0 & yes \\
GPT-4 over ONLINE-B & contested & no & 4 & 2 & no \\
GPT-4 over ONLINE-W & dominates & no & 0 & 0 & yes \\
GPT-4 over TranssionMT & contested & no & 3 & 2 & no \\
GPT-4 over Unbabel-Tower70B & contested & no & 4 & 2 & no \\
GPT-4 over refA & unresolved & no & 2 & 0 & yes \\
GPT-4 over refB & contested & no & 3 & 1 & no \\
Gemini-1.5-Pro over IKUN & dominates & yes & 0 & 0 & yes \\
Gemini-1.5-Pro over IKUN-C & dominates & yes & 0 & 0 & yes \\
Gemini-1.5-Pro over IOL-Research & dominates & no & 0 & 0 & no \\
Gemini-1.5-Pro over Llama3-70B & dominates & yes & 0 & 0 & yes \\
Gemini-1.5-Pro over MSLC & dominates & yes & 0 & 0 & yes \\
Mistral-Large over Gemini-1.5-Pro & contested & no & 7 & 2 & no \\
Gemini-1.5-Pro over ONLINE-A & dominates & yes & 0 & 0 & yes \\
ONLINE-B over Gemini-1.5-Pro & unresolved & no & 6 & 0 & no \\
Gemini-1.5-Pro over ONLINE-W & dominates & no & 0 & 0 & no \\
TranssionMT over Gemini-1.5-Pro & unresolved & no & 6 & 0 & no \\
Unbabel-Tower70B over Gemini-1.5-Pro & unresolved & no & 2 & 0 & no \\
refA over Gemini-1.5-Pro & contested & no & 7 & 2 & no \\
refB over Gemini-1.5-Pro & unresolved & no & 6 & 0 & no \\
IKUN over IKUN-C & dominates & yes & 0 & 0 & yes \\
IOL-Research over IKUN & dominates & yes & 0 & 0 & yes \\
Llama3-70B over IKUN & unresolved & no & 2 & 0 & no \\
IKUN over MSLC & dominates & yes & 0 & 0 & yes \\
Mistral-Large over IKUN & dominates & yes & 0 & 0 & yes \\
ONLINE-A over IKUN & dominates & no & 0 & 0 & yes \\
ONLINE-B over IKUN & dominates & yes & 0 & 0 & yes \\
ONLINE-W over IKUN & dominates & yes & 0 & 0 & yes \\
TranssionMT over IKUN & dominates & yes & 0 & 0 & yes \\
Unbabel-Tower70B over IKUN & dominates & yes & 0 & 0 & yes \\
refA over IKUN & dominates & yes & 0 & 0 & yes \\
refB over IKUN & dominates & yes & 0 & 0 & yes \\
IOL-Research over IKUN-C & dominates & yes & 0 & 0 & yes \\
Llama3-70B over IKUN-C & dominates & no & 0 & 0 & yes \\
IKUN-C over MSLC & dominates & yes & 0 & 0 & yes \\
Mistral-Large over IKUN-C & dominates & yes & 0 & 0 & yes \\
ONLINE-A over IKUN-C & dominates & yes & 0 & 0 & yes \\
ONLINE-B over IKUN-C & dominates & yes & 0 & 0 & yes \\
ONLINE-W over IKUN-C & dominates & yes & 0 & 0 & yes \\
TranssionMT over IKUN-C & dominates & yes & 0 & 0 & yes \\
Unbabel-Tower70B over IKUN-C & dominates & yes & 0 & 0 & yes \\
refA over IKUN-C & dominates & yes & 0 & 0 & yes \\
refB over IKUN-C & dominates & yes & 0 & 0 & yes \\
IOL-Research over Llama3-70B & dominates & yes & 0 & 0 & yes \\
IOL-Research over MSLC & dominates & yes & 0 & 0 & yes \\
Mistral-Large over IOL-Research & dominates & no & 0 & 0 & yes \\
IOL-Research over ONLINE-A & dominates & yes & 0 & 0 & yes \\
ONLINE-B over IOL-Research & dominates & no & 0 & 0 & yes \\
ONLINE-W over IOL-Research & dominates & no & 0 & 0 & no \\
TranssionMT over IOL-Research & dominates & no & 0 & 0 & yes \\
Unbabel-Tower70B over IOL-Research & dominates & yes & 0 & 0 & yes \\
refA over IOL-Research & dominates & no & 0 & 0 & no \\
refB over IOL-Research & dominates & no & 0 & 0 & yes \\
Llama3-70B over MSLC & dominates & yes & 0 & 0 & yes \\
Mistral-Large over Llama3-70B & dominates & yes & 0 & 0 & yes \\
ONLINE-A over Llama3-70B & dominates & no & 0 & 0 & no \\
ONLINE-B over Llama3-70B & dominates & yes & 0 & 0 & yes \\
ONLINE-W over Llama3-70B & dominates & yes & 0 & 0 & yes \\
TranssionMT over Llama3-70B & dominates & yes & 0 & 0 & yes \\
Unbabel-Tower70B over Llama3-70B & dominates & yes & 0 & 0 & yes \\
refA over Llama3-70B & dominates & yes & 0 & 0 & yes \\
refB over Llama3-70B & dominates & yes & 0 & 0 & yes \\
Mistral-Large over MSLC & dominates & yes & 0 & 0 & yes \\
ONLINE-A over MSLC & dominates & yes & 0 & 0 & yes \\
ONLINE-B over MSLC & dominates & yes & 0 & 0 & yes \\
ONLINE-W over MSLC & dominates & yes & 0 & 0 & yes \\
TranssionMT over MSLC & dominates & yes & 0 & 0 & yes \\
Unbabel-Tower70B over MSLC & dominates & yes & 0 & 0 & yes \\
refA over MSLC & dominates & yes & 0 & 0 & yes \\
refB over MSLC & dominates & yes & 0 & 0 & yes \\
Mistral-Large over ONLINE-A & dominates & yes & 0 & 0 & yes \\
ONLINE-B over Mistral-Large & dominates & no & 0 & 0 & no \\
Mistral-Large over ONLINE-W & dominates & no & 0 & 0 & no \\
TranssionMT over Mistral-Large & dominates & no & 0 & 0 & no \\
Unbabel-Tower70B over Mistral-Large & dominates & no & 0 & 0 & no \\
Mistral-Large over refA & unresolved & no & 2 & 0 & no \\
refB over Mistral-Large & unresolved & no & 2 & 0 & no \\
ONLINE-B over ONLINE-A & dominates & yes & 0 & 0 & yes \\
ONLINE-W over ONLINE-A & dominates & yes & 0 & 0 & yes \\
TranssionMT over ONLINE-A & dominates & yes & 0 & 0 & yes \\
Unbabel-Tower70B over ONLINE-A & dominates & yes & 0 & 0 & yes \\
refA over ONLINE-A & dominates & yes & 0 & 0 & yes \\
refB over ONLINE-A & dominates & yes & 0 & 0 & yes \\
ONLINE-B over ONLINE-W & dominates & no & 0 & 0 & yes \\
TranssionMT over ONLINE-B & unresolved & no & 5 & 0 & no \\
Unbabel-Tower70B over ONLINE-B & dominates & no & 0 & 0 & no \\
ONLINE-B over refA & dominates & no & 0 & 0 & no \\
ONLINE-B over refB & unresolved & no & 2 & 0 & no \\
TranssionMT over ONLINE-W & dominates & no & 0 & 0 & yes \\
Unbabel-Tower70B over ONLINE-W & dominates & yes & 0 & 0 & yes \\
refA over ONLINE-W & unresolved & no & 2 & 0 & no \\
refB over ONLINE-W & dominates & no & 0 & 0 & no \\
Unbabel-Tower70B over TranssionMT & dominates & no & 0 & 0 & no \\
TranssionMT over refA & dominates & no & 0 & 0 & no \\
TranssionMT over refB & unresolved & no & 2 & 0 & no \\
Unbabel-Tower70B over refA & dominates & no & 0 & 0 & yes \\
Unbabel-Tower70B over refB & dominates & no & 0 & 0 & no \\
refB over refA & dominates & no & 0 & 0 & no \\
\bottomrule\end{tabular}\end{center}

\paragraph{translation, WMT24 ja-zh, 15 systems.} Published order at ('wmt', 'none'): Claude-3.5 $>$ refA $>$ GPT-4 $>$ DLUT-GTCOM $>$ Unbabel-Tower70B $>$ Gemini-1.5-Pro $>$ CommandR-plus $>$ IOL-Research $>$ Aya23 $>$ Llama3-70B $>$ Team-J $>$ NTTSU $>$ ONLINE-B $>$ IKUN-C $>$ MSLC. 559 items. Accuracy by system and cell, from which every quantity below follows:

\begin{center}\small
\begin{tabular}{lcccccccc}
\toprule
system & wmt / none & wmt / clipped & plain / none & plain / clipped & major-only / none & major-only / clipped & count / none & count / clipped \\
\midrule
Claude-3.5 & -1.224 & -1.188 & -1.234 & -1.195 & -0.163 & -0.163 & -0.583 & -0.583 \\
refA & -1.325 & -1.307 & -1.300 & -1.290 & -0.179 & -0.179 & -0.585 & -0.585 \\
GPT-4 & -1.456 & -1.456 & -1.487 & -1.487 & -0.190 & -0.190 & -0.728 & -0.728 \\
DLUT-GTCOM & -1.519 & -1.449 & -1.531 & -1.506 & -0.216 & -0.216 & -0.665 & -0.665 \\
Unbabel-Tower70B & -1.695 & -1.673 & -1.734 & -1.714 & -0.240 & -0.240 & -0.775 & -0.775 \\
Gemini-1.5-Pro & -1.780 & -1.601 & -1.512 & -1.479 & -0.207 & -0.207 & -0.682 & -0.682 \\
CommandR-plus & -1.916 & -1.826 & -1.978 & -1.887 & -0.286 & -0.286 & -0.834 & -0.834 \\
IOL-Research & -2.108 & -1.999 & -1.978 & -1.884 & -0.286 & -0.286 & -0.834 & -0.834 \\
Aya23 & -3.039 & -2.889 & -2.817 & -2.750 & -0.413 & -0.413 & -1.165 & -1.165 \\
Llama3-70B & -3.072 & -2.954 & -2.798 & -2.746 & -0.404 & -0.404 & -1.181 & -1.181 \\
Team-J & -3.919 & -3.744 & -3.537 & -3.433 & -0.585 & -0.585 & -1.197 & -1.197 \\
NTTSU & -4.348 & -3.505 & -3.374 & -3.268 & -0.531 & -0.531 & -1.249 & -1.249 \\
ONLINE-B & -5.283 & -4.784 & -4.798 & -4.528 & -0.798 & -0.798 & -1.606 & -1.606 \\
IKUN-C & -6.610 & -5.944 & -4.934 & -4.766 & -0.889 & -0.889 & -1.377 & -1.377 \\
MSLC & -9.206 & -7.226 & -6.676 & -6.379 & -1.166 & -1.166 & -2.011 & -2.011 \\
\bottomrule\end{tabular}\end{center}

Per-pair classification. A pair is \emph{contested} only when some cell puts it the other way round with an interval excluding zero; a pair that is neither dominating nor contested is one the benchmark does not resolve, which is a different failure and is not counted as a reversal.

\begin{center}\small
\begin{tabular}{llccccl}
\toprule
pair (as the published cell orders it) & verdict & sep.\ every cell & cells reversing & with an interval & own cell resolves \\
\midrule
Claude-3.5 over Aya23 & dominates & yes & 0 & 0 & yes \\
CommandR-plus over Aya23 & dominates & yes & 0 & 0 & yes \\
DLUT-GTCOM over Aya23 & dominates & yes & 0 & 0 & yes \\
GPT-4 over Aya23 & dominates & yes & 0 & 0 & yes \\
Gemini-1.5-Pro over Aya23 & dominates & yes & 0 & 0 & yes \\
Aya23 over IKUN-C & dominates & yes & 0 & 0 & yes \\
IOL-Research over Aya23 & dominates & yes & 0 & 0 & yes \\
Aya23 over Llama3-70B & unresolved & no & 4 & 0 & no \\
Aya23 over MSLC & dominates & yes & 0 & 0 & yes \\
Aya23 over NTTSU & dominates & no & 0 & 0 & yes \\
Aya23 over ONLINE-B & dominates & yes & 0 & 0 & yes \\
Aya23 over Team-J & dominates & no & 0 & 0 & yes \\
Unbabel-Tower70B over Aya23 & dominates & yes & 0 & 0 & yes \\
refA over Aya23 & dominates & yes & 0 & 0 & yes \\
Claude-3.5 over CommandR-plus & dominates & yes & 0 & 0 & yes \\
Claude-3.5 over DLUT-GTCOM & dominates & no & 0 & 0 & no \\
Claude-3.5 over GPT-4 & dominates & no & 0 & 0 & no \\
Claude-3.5 over Gemini-1.5-Pro & dominates & no & 0 & 0 & yes \\
Claude-3.5 over IKUN-C & dominates & yes & 0 & 0 & yes \\
Claude-3.5 over IOL-Research & dominates & yes & 0 & 0 & yes \\
Claude-3.5 over Llama3-70B & dominates & yes & 0 & 0 & yes \\
Claude-3.5 over MSLC & dominates & yes & 0 & 0 & yes \\
Claude-3.5 over NTTSU & dominates & yes & 0 & 0 & yes \\
Claude-3.5 over ONLINE-B & dominates & yes & 0 & 0 & yes \\
Claude-3.5 over Team-J & dominates & yes & 0 & 0 & yes \\
Claude-3.5 over Unbabel-Tower70B & dominates & yes & 0 & 0 & yes \\
Claude-3.5 over refA & dominates & no & 0 & 0 & no \\
DLUT-GTCOM over CommandR-plus & dominates & yes & 0 & 0 & yes \\
GPT-4 over CommandR-plus & dominates & yes & 0 & 0 & yes \\
Gemini-1.5-Pro over CommandR-plus & dominates & no & 0 & 0 & no \\
CommandR-plus over IKUN-C & dominates & yes & 0 & 0 & yes \\
CommandR-plus over IOL-Research & unresolved & no & 6 & 0 & no \\
CommandR-plus over Llama3-70B & dominates & yes & 0 & 0 & yes \\
CommandR-plus over MSLC & dominates & yes & 0 & 0 & yes \\
CommandR-plus over NTTSU & dominates & yes & 0 & 0 & yes \\
CommandR-plus over ONLINE-B & dominates & yes & 0 & 0 & yes \\
CommandR-plus over Team-J & dominates & yes & 0 & 0 & yes \\
Unbabel-Tower70B over CommandR-plus & dominates & no & 0 & 0 & no \\
refA over CommandR-plus & dominates & yes & 0 & 0 & yes \\
GPT-4 over DLUT-GTCOM & unresolved & no & 3 & 0 & no \\
DLUT-GTCOM over Gemini-1.5-Pro & unresolved & no & 4 & 0 & no \\
DLUT-GTCOM over IKUN-C & dominates & yes & 0 & 0 & yes \\
DLUT-GTCOM over IOL-Research & dominates & yes & 0 & 0 & yes \\
DLUT-GTCOM over Llama3-70B & dominates & yes & 0 & 0 & yes \\
DLUT-GTCOM over MSLC & dominates & yes & 0 & 0 & yes \\
DLUT-GTCOM over NTTSU & dominates & yes & 0 & 0 & yes \\
DLUT-GTCOM over ONLINE-B & dominates & yes & 0 & 0 & yes \\
DLUT-GTCOM over Team-J & dominates & yes & 0 & 0 & yes \\
DLUT-GTCOM over Unbabel-Tower70B & dominates & no & 0 & 0 & no \\
refA over DLUT-GTCOM & dominates & no & 0 & 0 & no \\
GPT-4 over Gemini-1.5-Pro & unresolved & no & 3 & 0 & no \\
GPT-4 over IKUN-C & dominates & yes & 0 & 0 & yes \\
GPT-4 over IOL-Research & dominates & yes & 0 & 0 & yes \\
GPT-4 over Llama3-70B & dominates & yes & 0 & 0 & yes \\
GPT-4 over MSLC & dominates & yes & 0 & 0 & yes \\
GPT-4 over NTTSU & dominates & yes & 0 & 0 & yes \\
GPT-4 over ONLINE-B & dominates & yes & 0 & 0 & yes \\
GPT-4 over Team-J & dominates & yes & 0 & 0 & yes \\
GPT-4 over Unbabel-Tower70B & dominates & no & 0 & 0 & no \\
refA over GPT-4 & dominates & no & 0 & 0 & no \\
Gemini-1.5-Pro over IKUN-C & dominates & yes & 0 & 0 & yes \\
Gemini-1.5-Pro over IOL-Research & dominates & no & 0 & 0 & no \\
Gemini-1.5-Pro over Llama3-70B & dominates & yes & 0 & 0 & yes \\
Gemini-1.5-Pro over MSLC & dominates & yes & 0 & 0 & yes \\
Gemini-1.5-Pro over NTTSU & dominates & yes & 0 & 0 & yes \\
Gemini-1.5-Pro over ONLINE-B & dominates & yes & 0 & 0 & yes \\
Gemini-1.5-Pro over Team-J & dominates & yes & 0 & 0 & yes \\
Unbabel-Tower70B over Gemini-1.5-Pro & unresolved & no & 7 & 0 & no \\
refA over Gemini-1.5-Pro & dominates & no & 0 & 0 & yes \\
IOL-Research over IKUN-C & dominates & yes & 0 & 0 & yes \\
Llama3-70B over IKUN-C & dominates & yes & 0 & 0 & yes \\
IKUN-C over MSLC & dominates & yes & 0 & 0 & yes \\
NTTSU over IKUN-C & dominates & no & 0 & 0 & yes \\
ONLINE-B over IKUN-C & contested & no & 2 & 2 & yes \\
Team-J over IKUN-C & dominates & yes & 0 & 0 & yes \\
Unbabel-Tower70B over IKUN-C & dominates & yes & 0 & 0 & yes \\
refA over IKUN-C & dominates & yes & 0 & 0 & yes \\
IOL-Research over Llama3-70B & dominates & yes & 0 & 0 & yes \\
IOL-Research over MSLC & dominates & yes & 0 & 0 & yes \\
IOL-Research over NTTSU & dominates & yes & 0 & 0 & yes \\
IOL-Research over ONLINE-B & dominates & yes & 0 & 0 & yes \\
IOL-Research over Team-J & dominates & yes & 0 & 0 & yes \\
Unbabel-Tower70B over IOL-Research & dominates & no & 0 & 0 & yes \\
refA over IOL-Research & dominates & yes & 0 & 0 & yes \\
Llama3-70B over MSLC & dominates & yes & 0 & 0 & yes \\
Llama3-70B over NTTSU & dominates & no & 0 & 0 & yes \\
Llama3-70B over ONLINE-B & dominates & yes & 0 & 0 & yes \\
Llama3-70B over Team-J & dominates & no & 0 & 0 & yes \\
Unbabel-Tower70B over Llama3-70B & dominates & yes & 0 & 0 & yes \\
refA over Llama3-70B & dominates & yes & 0 & 0 & yes \\
NTTSU over MSLC & dominates & yes & 0 & 0 & yes \\
ONLINE-B over MSLC & dominates & yes & 0 & 0 & yes \\
Team-J over MSLC & dominates & yes & 0 & 0 & yes \\
Unbabel-Tower70B over MSLC & dominates & yes & 0 & 0 & yes \\
refA over MSLC & dominates & yes & 0 & 0 & yes \\
NTTSU over ONLINE-B & dominates & no & 0 & 0 & no \\
Team-J over NTTSU & unresolved & no & 5 & 0 & no \\
Unbabel-Tower70B over NTTSU & dominates & yes & 0 & 0 & yes \\
refA over NTTSU & dominates & yes & 0 & 0 & yes \\
Team-J over ONLINE-B & dominates & yes & 0 & 0 & yes \\
Unbabel-Tower70B over ONLINE-B & dominates & yes & 0 & 0 & yes \\
refA over ONLINE-B & dominates & yes & 0 & 0 & yes \\
Unbabel-Tower70B over Team-J & dominates & yes & 0 & 0 & yes \\
refA over Team-J & dominates & yes & 0 & 0 & yes \\
refA over Unbabel-Tower70B & dominates & yes & 0 & 0 & yes \\
\bottomrule\end{tabular}\end{center}

\paragraph{docking, 7 programs.} Published order at ('rmsd2+valid', 'none'): vina $>$ gold $>$ diffdock $>$ deepdock $>$ tankbind $>$ unimol $>$ equibind. 428 items. Accuracy by system and cell, from which every quantity below follows:

\begin{center}\small
\begin{tabular}{lcccccccc}
\toprule
system & rmsd2 / none & rmsd2 / energy minimization & rmsd2+valid / none & rmsd2+valid / energy minimization & rmsd1 / none & rmsd1 / energy minimization & rmsd2+chem / none & rmsd2+chem / energy minimization \\
\midrule
vina & 0.523 & 0.484 & 0.512 & 0.421 & 0.344 & 0.306 & 0.519 & 0.458 \\
gold & 0.512 & 0.479 & 0.477 & 0.423 & 0.344 & 0.292 & 0.509 & 0.470 \\
diffdock & 0.379 & 0.381 & 0.136 & 0.334 & 0.145 & 0.187 & 0.355 & 0.357 \\
deepdock & 0.178 & 0.196 & 0.047 & 0.129 & 0.047 & 0.072 & 0.133 & 0.154 \\
tankbind & 0.149 & 0.199 & 0.026 & 0.119 & 0.026 & 0.072 & 0.072 & 0.133 \\
unimol & 0.229 & 0.280 & 0.019 & 0.175 & 0.042 & 0.101 & 0.063 & 0.194 \\
equibind & 0.026 & 0.054 & 0.002 & 0.047 & 0.000 & 0.023 & 0.026 & 0.051 \\
\bottomrule\end{tabular}\end{center}

Per-pair classification. A pair is \emph{contested} only when some cell puts it the other way round with an interval excluding zero; a pair that is neither dominating nor contested is one the benchmark does not resolve, which is a different failure and is not counted as a reversal.

\begin{center}\small
\begin{tabular}{llccccl}
\toprule
pair (as the published cell orders it) & verdict & sep.\ every cell & cells reversing & with an interval & own cell resolves \\
\midrule
diffdock over deepdock & dominates & yes & 0 & 0 & yes \\
deepdock over equibind & dominates & yes & 0 & 0 & yes \\
gold over deepdock & dominates & yes & 0 & 0 & yes \\
deepdock over tankbind & unresolved & no & 2 & 0 & no \\
deepdock over unimol & contested & no & 5 & 2 & yes \\
vina over deepdock & dominates & yes & 0 & 0 & yes \\
diffdock over equibind & dominates & yes & 0 & 0 & yes \\
gold over diffdock & dominates & yes & 0 & 0 & yes \\
diffdock over tankbind & dominates & yes & 0 & 0 & yes \\
diffdock over unimol & dominates & yes & 0 & 0 & yes \\
vina over diffdock & dominates & yes & 0 & 0 & yes \\
gold over equibind & dominates & yes & 0 & 0 & yes \\
tankbind over equibind & dominates & yes & 0 & 0 & yes \\
unimol over equibind & dominates & yes & 0 & 0 & yes \\
vina over equibind & dominates & yes & 0 & 0 & yes \\
gold over tankbind & dominates & yes & 0 & 0 & yes \\
gold over unimol & dominates & yes & 0 & 0 & yes \\
vina over gold & unresolved & no & 3 & 0 & no \\
tankbind over unimol & contested & no & 6 & 4 & no \\
vina over tankbind & dominates & yes & 0 & 0 & yes \\
vina over unimol & dominates & yes & 0 & 0 & yes \\
\bottomrule\end{tabular}\end{center}

\paragraph{The nine \textsc{esa} translation boards.} Their per-cell accuracy matrices are not printed: nine boards of eleven to sixteen systems by eight cells is four hundred numbers that no reader checks, and they are in \texttt{results/robust\_order\_wmt24\_esa.json} in full. What the text uses from them is here.

\begin{center}\footnotesize

\begin{tabular}{lccccccc}

\toprule

pair & systems & segments & annotators & dominate & contested (cert.) & unresolved & tiers \\

\midrule

cs-uk & 12 & 1954 & 14 & 58 of 66 & 2 (1) & 6 & 7 \\

en-cs & 16 & 634 & 19 & 101 of 120 & 4 (0) & 15 & 8 \\

en-es & 14 & 634 & 14 & 79 of 91 & 4 (1) & 8 & 9 \\

en-hi & 11 & 634 & 15 & 39 of 55 & 7 (2) & 9 & 7 \\

en-is & 11 & 634 & 4 & 50 of 55 & 2 (1) & 3 & 8 \\

en-ja & 13 & 634 & 14 & 60 of 78 & 6 (0) & 12 & 7 \\

en-ru & 14 & 634 & 7 & 77 of 91 & 3 (2) & 11 & 8 \\

en-uk & 11 & 634 & 8 & 45 of 55 & 0 (0) & 10 & 7 \\

en-zh & 13 & 634 & 14 & 50 of 78 & 9 (1) & 19 & 5 \\

\bottomrule

\end{tabular}

\end{center}



\paragraph{The eight boards of the previous edition.} Same treatment, and the column a reader wants that the later edition cannot supply is the last but one: how many ratings a segment carries, which is what gives this protocol its criterion axis.

\begin{center}\footnotesize

\begin{tabular}{lccccccc c}

\toprule

pair & systems & segments & annotators & ratings & dominate & contested (cert.) & unresolved & tiers \\

\midrule

ces-ukr & 14 & 1125 & 146 & 1.792 & 81 of 91 & 4 (0) & 6 & 10 \\

deu-eng & 14 & 716 & 82 & 2.037 & 86 of 91 & 1 (0) & 4 & 10 \\

eng-ces & 16 & 1238 & 162 & 1.581 & 98 of 120 & 5 (2) & 17 & 8 \\

eng-deu & 13 & 741 & 87 & 2.127 & 66 of 78 & 2 (0) & 10 & 8 \\

eng-jpn & 17 & 1238 & 164 & 1.52 & 124 of 136 & 3 (0) & 9 & 10 \\

eng-zho & 16 & 1238 & 154 & 1.52 & 107 of 120 & 0 (0) & 13 & 11 \\

jpn-eng & 18 & 1302 & 174 & 1.495 & 141 of 153 & 1 (0) & 11 & 13 \\

zho-eng & 16 & 1015 & 128 & 1.569 & 105 of 120 & 8 (1) & 7 & 10 \\

\bottomrule

\end{tabular}

\end{center}



